% Preamble template for Tufte-style applied econometrics lecture notes.
% Do not compile this file directly; compile the main notes file instead.
\documentclass{tufte-handout}

\usepackage[T1]{fontenc}
\usepackage[utf8]{inputenc}
\usepackage{ntheorem}
\usepackage{graphicx}
\usepackage{amsmath}
\usepackage{amssymb}
\usepackage{mathtools}
\usepackage{hyperref}
\usepackage{epigraph}
\usepackage{booktabs}
\usepackage{array}
\theoremstyle{break}

\hypersetup{
  colorlinks,
  urlcolor = blue,
  pdfauthor={Swapnil Singh},
  pdfkeywords={econometrics, causal inference, applied econometrics},
  pdftitle={Lecture Notes for Applied Econometrics},
  pdfpagemode=UseNone
}

\newtheorem{ruleN}{Rule}
\newtheorem{thmN}{Theorem}
\newtheorem{assN}{Assumption}
\newtheorem{defN}{Definition}
\newtheorem{exmp}{Example}
\newtheorem{cmt}{Comment}
\newtheorem{proof}{Proof}

\newcommand{\continuation}{??}
\newtheorem*{excont}{Example \continuation}
\newenvironment{continueexample}[1]
 {\renewcommand{\continuation}{\ref{#1}}\excont[continued]}
 {\endexcont}

\newcommand{\bY}{\mathbf{Y}}
\newcommand{\bX}{\mathbf{X}}
\newcommand{\bD}{\mathbf{D}}
\newcommand{\E}{\mathbb{E}}
\newcommand{\Prob}{\mathbb{P}}
\newcommand{\Var}{\mathrm{Var}}
\newcommand{\Cov}{\mathrm{Cov}}
\newcommand{\se}{\mathrm{SE}}
\newcommand{\CI}{\mathrm{CI}}
\newcommand{\ATE}{\mathrm{ATE}}
\newcommand{\ATT}{\mathrm{ATT}}
\newcommand{\TOT}{\mathrm{TOT}}
\newcommand{\LATE}{\mathrm{LATE}}
\newcommand{\ITT}{\mathrm{ITT}}
\newcommand{\RD}{\mathrm{RD}}
\newcommand{\DiD}{\mathrm{DiD}}
\newcommand{\MMSE}{\mathrm{MMSE}}
\newcommand{\OLS}{\mathrm{OLS}}
\newcommand{\Normal}{\mathcal{N}}
\newcommand{\plim}{\operatorname*{plim}}
\DeclareMathOperator*{\argmin}{arg\,min}
\DeclareMathOperator{\Supp}{Supp}

\newcommand\independent{\protect\mathpalette{\protect\independenT}{\perp}}
\def\independenT#1#2{\mathrel{\rlap{$#1#2$}\mkern2mu{#1#2}}}
\newcommand{\indep}{\perp\!\!\!\perp}

\usepackage{cleveref}
\crefname{appsec}{appendix}{appendices}
\crefname{appsubsec}{appendix}{appendices}
\crefname{assumption}{assumption}{assumptions}
\crefname{equation}{equation}{equations}
\crefname{exmp}{example}{examples}
\crefname{assN}{assumption}{assumptions}
\crefname{cmt}{comment}{comments}
\crefname{defN}{definition}{definitions}

\usepackage[nolist]{acronym}
\begin{acronym}
  \acro{ATE}{average treatment effect}
  \acro{ATT}{average treatment effect on the treated}
  \acro{CIA}{conditional independence assumption}
  \acro{CI}{confidence interval}
  \acro{CLT}{central limit theorem}
  \acro{DiD}{difference-in-differences}
  \acro{FE}{fixed effects}
  \acro{FWL}{Frisch-Waugh-Lovell}
  \acro{IV}{instrumental variables}
  \acro{LATE}{local average treatment effect}
  \acro{OLS}{ordinary least squares}
  \acro{OVB}{omitted-variable bias}
  \acro{PO}{potential outcomes}
  \acro{RCT}{randomized controlled trial}
  \acro{RD}{regression discontinuity}
  \acro{SUTVA}{stable unit treatment value assignment}
  \acro{TWFE}{two-way fixed effects}
\end{acronym}

\def\inprobHIGH{\,{\buildrel p \over \rightarrow}\,}
\def\inprob{\,{\inprobHIGH}\,}
\def\indistHIGH{\,{\buildrel d \over \rightarrow}\,}
\def\indist{\,{\indistHIGH}\,}

\usepackage[many]{tcolorbox}
\definecolor{main}{HTML}{3F6EA7}
\definecolor{sub}{HTML}{F2F7FD}
\definecolor{sub2}{HTML}{FFF5E8}
\definecolor{mainline}{HTML}{D3DFEE}
\definecolor{warn}{HTML}{C7771B}
\definecolor{warnline}{HTML}{EAD9C0}

\tcbset{
    enhanced,
    breakable,
    colback = white,
    boxrule = 0.5pt,
    arc = 2pt,
    left = 9pt,
    right = 9pt,
    top = 7pt,
    bottom = 7pt,
    before skip = 0.3cm,
    after skip = 0.45cm
}

\newtcolorbox{boxD}{
    colback = sub,
    colframe = mainline,
    borderline west = {2.2pt}{0pt}{main},
    borderline north = {0.6pt}{0pt}{main!25},
    borderline south = {0.6pt}{0pt}{main!15}
}

\newtcolorbox{boxF}{
    colback = sub2,
    colframe = warnline,
    borderline west = {2.2pt}{0pt}{warn},
    borderline north = {0.6pt}{0pt}{warn!25},
    borderline south = {0.6pt}{0pt}{warn!15}
}

\usepackage{colortbl}
\usepackage{tikz}
\usepackage{verbatim}
\usetikzlibrary{positioning}
\usetikzlibrary{snakes}
\usetikzlibrary{calc}
\usetikzlibrary{arrows}
\usetikzlibrary{decorations.markings}
\usetikzlibrary{shapes.misc}
\usetikzlibrary{matrix,shapes,arrows,fit,tikzmark}


\title{Applied Econometrics: Session 6\\Matching and Propensity Scores}
\author{Swapnil Singh}
\date{}

\hypersetup{
  pdftitle={Applied Econometrics: Session 6 - Matching and Propensity Scores},
  pdfauthor={Swapnil Singh},
  pdfkeywords={matching, propensity score, conditional independence, overlap, inverse probability weighting}
}

\setlength{\headheight}{26pt}

\newcommand{\ATU}{\mathrm{ATU}}
\newcommand{\ps}{e(X)}
\newcommand{\ehat}{\widehat e(X)}

\begin{document}
\maketitle

\begin{abstract}
This lecture studies matching and propensity-score methods for observational
causal inference. The main question is simple: when treatment is not randomly
assigned, can we construct a credible comparison group using observed
pre-treatment covariates? We cover conditional independence, overlap, ATE and
ATT, exact matching, nearest-neighbor matching, common support, the
propensity-score theorem, propensity-score matching, inverse probability
weighting, diagnostics, and reporting standards.
\end{abstract}

\begin{boxD}
\textbf{Reading and objective.} Read Chapters 10 and 11 of Matheus Facure
Alves, \emph{Causal Inference for the Brave and True}:
\url{https://matheusfacure.github.io/python-causality-handbook/10-Matching.html}
and
\url{https://matheusfacure.github.io/python-causality-handbook/11-Propensity-Score.html}.
The companion script is
\texttt{lectures/code/lecture-6-matching-propensity.py}. The objective is to
learn how to design, diagnose, estimate, and report a selection-on-observables
study without pretending that software has solved the identification problem.
\end{boxD}

\section{The Observational Problem}

We observe a sample
\[
    (Y_i,D_i,X_i),\qquad D_i\in\{0,1\},
\]
where \(Y_i\) is the observed outcome, \(D_i\) is the treatment indicator, and
\(X_i\) is a vector of pre-treatment covariates. The potential outcomes are
\(Y_{1i}\), the outcome unit \(i\) would have under treatment, and \(Y_{0i}\),
the outcome unit \(i\) would have without treatment. The observed outcome is
\[
    Y_i = D_iY_{1i} + (1-D_i)Y_{0i}.
\]
For each unit we observe only one potential outcome. If \(D_i=1\), then
\(Y_i=Y_{1i}\) and \(Y_{0i}\) is missing. If \(D_i=0\), then \(Y_i=Y_{0i}\) and
\(Y_{1i}\) is missing. This is the fundamental missing-data problem of causal
inference.

In a randomized experiment, treated and untreated units are comparable in large
samples because assignment is independent of potential outcomes. In an
observational study, treatment is often chosen by people, firms, doctors,
schools, or governments. Treated and untreated units can differ even before the
treatment occurs. Those pre-treatment differences may also predict outcomes.
Then the raw treated-control difference,
\[
    \E[Y_i\mid D_i=1]-\E[Y_i\mid D_i=0],
\]
mixes the causal effect with selection bias.

\begin{defN}[Selection bias]
Selection bias is the difference between the observed comparison group and the
missing counterfactual group. For the treated population,
\[
\begin{aligned}
    &\E[Y_i\mid D_i=1]-\E[Y_i\mid D_i=0] \\
    &\quad =
    \ATT
    +
    \{\E[Y_{0i}\mid D_i=1]
    -\E[Y_{0i}\mid D_i=0]\}.
\end{aligned}
\]
The second term is selection bias. It is zero in a randomized experiment, but
not generally zero in observational data.
\end{defN}

\begin{exmp}[Training program]
Suppose participants in a job-training program earn more after the program than
non-participants. Is that the effect of training? Maybe. But participants may
also be younger, more motivated, closer to employment offices, or more likely
to have recent employment histories. If those characteristics would have raised
earnings even without training, the naive gap overstates the effect. Matching
asks whether we can compare trained and untrained people with similar
pre-treatment characteristics.
\end{exmp}

\section{Identification Assumptions}

\begin{assN}[Conditional independence]
The conditional independence assumption, also called selection on observables or
unconfoundedness, is
\[
    (Y_{1i},Y_{0i}) \independent D_i \mid X_i.
\]
After conditioning on \(X_i\), treatment assignment is as good as random.
\end{assN}

This assumption is not a mechanical assumption. It is a claim about the economic
assignment process. If the treatment is a scholarship, \(X_i\) must contain the
pre-treatment determinants of both scholarship receipt and later outcomes:
prior grades, family background, school characteristics, location, and so on.
If unobserved motivation affects both scholarship receipt and later income,
conditioning on observed \(X_i\) will not remove that bias.

\begin{boxF}
\textbf{Important warning.} Matching and propensity scores do not fix selection
on unobservables. They organize comparisons among units that look similar in
observed variables. If the important difference is unobserved, these methods
can still be badly biased.
\end{boxF}

\begin{assN}[Overlap]
For the covariate values relevant to the target estimand,
\[
    0 < \Prob(D_i=1\mid X_i=x) < 1.
\]
There must be treated and untreated units with comparable \(X_i\).
\end{assN}

Overlap is the empirical content behind the phrase ``comparable group.'' If all
high-income urban workers receive training and no high-income urban worker is
untreated, the data contain no direct evidence about what those treated workers
would have earned without training. An estimator can extrapolate, but that is
modeling, not matching.

\subsection{ATE, ATT, and ATU}

Three estimands appear often:
\[
    \ATE = \E[Y_{1i}-Y_{0i}],
\]
\[
    \ATT = \E[Y_{1i}-Y_{0i}\mid D_i=1],
\]
\[
    \ATU = \E[Y_{1i}-Y_{0i}\mid D_i=0].
\]
The ATE averages over the whole population. The ATT averages over the treated
population. The ATU averages over the untreated population. If treatment effects
are heterogeneous, these are different objects.

Many matching exercises naturally estimate the ATT because we start with each
treated unit and search for untreated matches to impute its missing \(Y_0\).
IPW can be written to target either the ATE or the ATT. Therefore the estimand
must be chosen before estimation. Do not report ``the effect'' without saying
whose effect.

\section{Exact Matching}

Exact matching is the cleanest version of the idea. Suppose \(X_i\) takes a
small number of discrete values. For a covariate cell \(x\), define the within
cell treatment-control difference
\[
    \delta(x)
    =
    \E[Y_i\mid D_i=1,X_i=x]
    -
    \E[Y_i\mid D_i=0,X_i=x].
\]
Under conditional independence and overlap,
\[
    \delta(x)
    =
    \E[Y_{1i}-Y_{0i}\mid X_i=x].
\]
For the ATT,
\[
    \ATT
    =
    \sum_x \delta(x)\Prob(X_i=x\mid D_i=1).
\]
The weights are the distribution of \(X_i\) among treated units, because the
target population is the treated population.

\begin{exmp}[A tiny exact-matching calculation]
Suppose four covariate cells produce the following cell effects and treated
shares.
\[
\begin{array}{lccc}
\toprule
\text{Cell} & \text{Cell effect} & \text{Treated share} & \text{Contribution}\\
\midrule
\text{Young, urban} & 2.5 & 0.40 & 1.00\\
\text{Young, rural} & 2.0 & 0.20 & 0.40\\
\text{Older, urban} & 1.5 & 0.25 & 0.375\\
\text{Older, rural} & 1.0 & 0.15 & 0.150\\
\bottomrule
\end{array}
\]
The exact-matching ATT is
\[
    1.00+0.40+0.375+0.150=1.925.
\]
Notice that we do not compare all treated units with all controls. We compare
within cells, then average using treated-population weights.
\end{exmp}

\subsection{Exact Matching as the G-Formula}

Start from
\[
    \ATT = \E[Y_{1i}-Y_{0i}\mid D_i=1].
\]
Iterate over covariates:
\[
\begin{aligned}
    \ATT
    =
    \sum_x
    \left\{
    \E[Y_{1i}\mid X_i=x,D_i=1]\right.\\
    \left.
    -
    \E[Y_{0i}\mid X_i=x,D_i=1]
    \right\}
    \Prob(X_i=x\mid D_i=1).
\end{aligned}
\]
The first conditional mean is observed:
\[
    \E[Y_{1i}\mid X_i=x,D_i=1]=\E[Y_i\mid X_i=x,D_i=1].
\]
For the missing mean, conditional independence gives
\[
\begin{aligned}
    \E[Y_{0i}\mid X_i=x,D_i=1]
    &=
    \E[Y_{0i}\mid X_i=x,D_i=0]\\
    &=
    \E[Y_i\mid X_i=x,D_i=0].
\end{aligned}
\]
Substituting yields the exact-matching formula above. This derivation is useful
because it shows precisely where the identifying assumption enters.

\section{Matching as Imputation}

Exact matching becomes difficult when \(X_i\) is continuous or high
dimensional. Nearest-neighbor matching replaces exact equality with closeness.
For a treated unit \(i\), let \(\mathcal{N}_K(i;D=0)\) be the set of its
\(K\) nearest untreated neighbors. Then
\[
    \widehat Y_{0i}
    =
    \frac{1}{K}\sum_{\ell\in\mathcal{N}_K(i;D=0)}Y_\ell,
\]
and the matching estimator for the ATT is
\[
    \widehat{\ATT}_M
    =
    \frac{1}{N_1}\sum_{i:D_i=1}(Y_i-\widehat Y_{0i}).
\]
This is imputation. We fill in the missing untreated potential outcome for each
treated unit using observed untreated outcomes from similar units.

\begin{defN}[Distance metric]
A common distance between treated unit \(i\) and control unit \(\ell\) is
\[
    d(i,\ell)
    =
    \left[(X_i-X_\ell)'W(X_i-X_\ell)\right]^{1/2}.
\]
If \(W=I\), this is Euclidean distance after scaling. If
\(W=\widehat\Sigma_X^{-1}\), this is Mahalanobis distance.
\end{defN}

Scaling is not optional. If one variable is income in euros and another is age
in years, raw Euclidean distance will usually be dominated by income. Standard
practice is to standardize continuous covariates before distance matching.

\subsection{Replacement, Calipers, and Neighbors}

Matching with replacement allows the same control unit to be used many times.
This often improves match quality because each treated unit can use the closest
available control. The cost is that the estimator may depend heavily on a small
number of controls.

Matching without replacement uses each control at most once. This can increase
the number of distinct controls, but later treated units may receive poor
matches because the best controls have already been used.

A caliper imposes a maximum acceptable distance:
\[
    d(i,\ell) \leq c.
\]
Calipers prevent visibly bad matches. The cost is that some treated units are
dropped, so the estimand may become the ATT for the matched treated population,
not the ATT for all treated units.

\begin{boxD}
\textbf{Bias-variance tradeoff.} One close neighbor is local and often has low
bias, but it is noisy. Many neighbors reduce variance but can introduce bias by
including less comparable controls. Strict calipers improve credibility of
matches but may change the target population.
\end{boxD}

\subsection{Bias from Imperfect Matches}

Let
\[
    m_0(x)=\E[Y_{0i}\mid X_i=x].
\]
Suppose treated unit \(i\) is matched to control \(m(i)\). Conditional on their
covariates, the expected imputation error is approximately
\[
    \E[Y_{0,m(i)}-Y_{0i}\mid X_i,X_{m(i)}]
    =
    m_0(X_{m(i)})-m_0(X_i).
\]
If \(m_0(\cdot)\) is smooth, a first-order Taylor approximation gives
\[
    m_0(X_{m(i)})-m_0(X_i)
    \approx
    \nabla m_0(X_i)'(X_{m(i)}-X_i).
\]
The bias is small when matched covariates are close or when the untreated
outcome surface is flat. It is large when the match is far away in a direction
that strongly predicts untreated outcomes.

\subsection{Curse of Dimensionality}

Suppose each covariate is split into five bins. With \(p\) covariates there are
\[
    5^p
\]
cells. With \(p=2\), there are 25 cells. With \(p=5\), there are 3,125 cells.
With \(p=10\), there are 9,765,625 cells. Most cells will be empty in ordinary
datasets. Even nearest-neighbor matching suffers: in high dimensions, the
nearest neighbor may still be far away.

This is why matching is attractive as an idea but difficult as a mechanical
solution. Adding covariates can reduce omitted confounding, but it also makes
good matches harder to find. The answer is not to throw in every available
variable blindly. The answer is to choose pre-treatment confounders using
economic reasoning, then diagnose whether the data contain credible matches.

\section{Common Support and Design Diagnostics}

Common support is not a decoration added after estimation. It is part of the
research design. Ask:
\begin{itemize}
    \item Do treated units have nearby controls?
    \item Do controls have nearby treated units if the target is the ATE?
    \item How many observations are dropped by calipers or trimming?
    \item Are results driven by a small set of heavily reused controls?
    \item Are covariates balanced after matching?
\end{itemize}

A matched sample can have fewer observations than the original sample. This is
not automatically a problem. A smaller credible comparison is often better than
a larger incredible comparison. But the population must be described honestly.
If the method drops older high-income treated units because no controls are
similar, the estimate is not the ATT for all treated units.

\section{Propensity Scores}

\begin{defN}[Propensity score]
The propensity score is the conditional probability of treatment:
\[
    e(X_i)=\Prob(D_i=1\mid X_i).
\]
\end{defN}

Propensity scores compress the treatment-selection information in \(X_i\) into
a scalar. Units with the same propensity score have the same predicted
probability of treatment, even if their covariate profiles are not identical.
This can make matching, stratification, trimming, and weighting more feasible.

\begin{thmN}[Propensity-score theorem]
If
\[
    (Y_{1i},Y_{0i})\independent D_i\mid X_i,
\]
then
\[
    (Y_{1i},Y_{0i})\independent D_i\mid e(X_i),
\]
provided overlap holds. The propensity score is a balancing score.
\end{thmN}

\subsection{Balancing Intuition}

The key balancing statement is
\[
    D_i \independent X_i \mid e(X_i).
\]
Within a narrow propensity-score stratum, treated and untreated units should
have similar distributions of the covariates used to estimate the score. This
does not mean each treated unit has a perfect twin. It means the distribution of
covariates should be balanced between treatment groups within score strata or
after weighting.

\begin{proof}[Why the score balances covariates]
Let \(e(X)=\Prob(D=1\mid X)\). For any value \(p\) in the support of \(e(X)\),
\[
    \Prob(D=1\mid X=x,e(X)=p)=\Prob(D=1\mid X=x)=e(x)=p.
\]
Therefore, within the set of units with \(e(X)=p\), the probability of
treatment does not depend on the remaining details of \(X\). This is the sense
in which the score balances covariates.
\end{proof}

\begin{proof}[Why conditioning on the score identifies effects]
If conditional independence holds given \(X\), then for \(d\in\{0,1\}\),
\[
    \E[Y_d\mid D=1,e(X)]
    =
    \E\{\E[Y_d\mid D=1,X]\mid D=1,e(X)\}.
\]
By conditional independence,
\[
    \E[Y_d\mid D=1,X]=\E[Y_d\mid X].
\]
By the balancing property, the distribution of \(X\) conditional on
\(D=1,e(X)\) equals the distribution conditional on \(D=0,e(X)\). Therefore
\[
    \E[Y_d\mid D=1,e(X)] = \E[Y_d\mid D=0,e(X)].
\]
Thus treated and untreated units with the same propensity score can be compared
under the same selection-on-observables assumption.
\end{proof}

\begin{boxF}
\textbf{Beginner trap.} A propensity-score model is not judged mainly by its
classification accuracy or AUC. A model that predicts treatment almost
perfectly may create terrible overlap. The goal is balance and common support,
not a beautiful treatment-choice regression table.
\end{boxF}

\section{Estimating Propensity Scores}

In practice \(e(X_i)\) is unknown, so we estimate it:
\[
    \widehat e(X_i)=\widehat{\Prob}(D_i=1\mid X_i).
\]
A common starting point is logistic regression,
\[
    \log\left(\frac{e(X_i)}{1-e(X_i)}\right)=X_i'\gamma.
\]
Researchers may also use probit, series terms, interactions, random forests, or
gradient boosting. More flexible models can help balance, but they can also
overfit. In applied work, the score model should be treated as a design model:
its job is to make treated and untreated units comparable on pre-treatment
covariates.

\begin{boxD}
\textbf{Covariate choice.} Include pre-treatment variables that predict both
treatment and potential outcomes. Do not include variables caused by treatment.
Be careful with pure instruments: variables that strongly predict treatment but
not outcomes can worsen overlap and inflate weights.
\end{boxD}

\subsection{Diagnostics After Estimating the Score}

After estimating \(\widehat e(X)\), inspect:
\begin{itemize}
    \item histograms or density plots of scores by treatment status;
    \item minimum, median, and maximum score in each group;
    \item the fraction of observations near 0 or 1;
    \item standardized mean differences before and after adjustment;
    \item maximum weights and effective sample size;
    \item which units are dropped by common-support restrictions.
\end{itemize}

For a covariate \(X_j\), the standardized mean difference is
\[
    \mathrm{SMD}_j
    =
    \frac{\bar X_{j,1}-\bar X_{j,0}}{s_j},
\]
where \(s_j\) is a pooled or full-sample standard deviation. Weighted versions
replace the group means by weighted means. Values close to zero indicate better
balance. The threshold 0.10 is often used as a rule of thumb, but it is not a
law of nature.

\section{Trimming and Propensity-Score Matching}

Trimming removes observations with weak overlap:
\[
    a \leq \widehat e(X_i) \leq b.
\]
Common choices include fixed cutoffs such as \(0.05\leq \widehat e(X_i)\leq
0.95\), or the empirical common support. Define
\[
    \ell
    =
    \max\{\min_{D_i=1}\widehat e(X_i),\min_{D_i=0}\widehat e(X_i)\},
    \qquad
    u
    =
    \min\{\max_{D_i=1}\widehat e(X_i),\max_{D_i=0}\widehat e(X_i)\}.
\]
Then keep observations satisfying
\[
    \ell \leq \widehat e(X_i)\leq u.
\]
Trimming is not a harmless numerical step. It changes the estimand to the
population that remains after trimming.

Propensity-score matching uses the scalar distance
\[
    d(i,\ell)=|\widehat e(X_i)-\widehat e(X_\ell)|.
\]
The advantage is simplicity. The danger is that two units with the same score
can have different covariate profiles. Therefore balance must be checked on the
original covariates, not only on the score.

\section{Inverse Probability Weighting}

Inverse probability weighting reweights observations by how rare their
treatment status is for their covariates. For the ATE, the population moments
are identified by
\[
    \E[Y_1] = \E\left[\frac{D_iY_i}{e(X_i)}\right],
    \qquad
    \E[Y_0] = \E\left[\frac{(1-D_i)Y_i}{1-e(X_i)}\right].
\]
Therefore
\[
    \ATE
    =
    \E\left[\frac{D_iY_i}{e(X_i)}
    -
    \frac{(1-D_i)Y_i}{1-e(X_i)}\right].
\]

\begin{proof}[IPW identity for \(Y_1\)]
Using iterated expectations,
\[
    \E\left[\frac{D_iY_i}{e(X_i)}\right]
    =
    \E\left[\E\left(\frac{D_iY_{1i}}{e(X_i)}\mid X_i\right)\right].
\]
Under conditional independence,
\[
    \E[D_iY_{1i}\mid X_i]
    =
    \E[D_i\mid X_i]\E[Y_{1i}\mid X_i]
    =
    e(X_i)\E[Y_{1i}\mid X_i].
\]
Thus
\[
    \E\left[\frac{D_iY_i}{e(X_i)}\right]
    =
    \E[\E(Y_{1i}\mid X_i)]
    =
    \E[Y_{1i}].
\]
The proof for \(Y_0\) is analogous.
\end{proof}

A normalized, Hajek-style ATE estimator is
\[
    \widehat{\ATE}_{IPW}
    =
    \frac{\sum_i D_iY_i/\widehat e(X_i)}
         {\sum_i D_i/\widehat e(X_i)}
    -
    \frac{\sum_i (1-D_i)Y_i/[1-\widehat e(X_i)]}
         {\sum_i (1-D_i)/[1-\widehat e(X_i)]}.
\]
The treated group is weighted to represent everyone. The untreated group is
also weighted to represent everyone. The difference estimates the ATE under
conditional independence and overlap.

\subsection{IPW for the ATT}

For the ATT, treated units already represent the treated population. Controls
must be reweighted to look like treated units:
\[
    w_i^{ATT}
    =
    \begin{cases}
      1, & D_i=1,\\
      \dfrac{e(X_i)}{1-e(X_i)}, & D_i=0.
    \end{cases}
\]
A normalized ATT estimator is
\[
    \widehat{\ATT}_{IPW}
    =
    \frac{\sum_i D_iY_i}{\sum_iD_i}
    -
    \frac{\sum_i (1-D_i)w_i^{ATT}Y_i}
         {\sum_i(1-D_i)w_i^{ATT}}.
\]
Controls with high treatment probability receive larger weight because they
look like the treated population.

\subsection{Weight Diagnostics}

Extreme propensity scores create extreme weights. For the ATE, a treated unit
with \(\widehat e(X_i)\) close to zero receives weight \(1/\widehat e(X_i)\).
A control unit with \(\widehat e(X_i)\) close to one receives weight
\(1/[1-\widehat e(X_i)]\). A few observations can dominate the estimate.

Report at least the maximum weight, selected percentiles, and effective sample
size:
\[
    ESS = \frac{(\sum_i w_i)^2}{\sum_i w_i^2}.
\]
If 2,000 observations have an effective sample size of 60 after weighting, the
design is fragile even if the regression output prints a small standard error.

\section{The Companion Python Example}

The script \texttt{lectures/code/lecture-6-matching-propensity.py} simulates an
observational study with selection on observables. Treatment is more likely for
some observed covariate profiles, treatment effects are heterogeneous, and the
true ATE and ATT are known because the data are simulated.

Students should run
\[
    \texttt{python lectures/code/lecture-6-matching-propensity.py}
\]
and compare:
\begin{itemize}
    \item the naive treated-control gap;
    \item regression-adjusted OLS;
    \item one-nearest-neighbor matching for the ATT;
    \item propensity-score IPW for the ATE;
    \item propensity-score weights for the ATT;
    \item trimmed IPW estimates.
\end{itemize}
The useful lesson is not that one estimator always wins. The useful lesson is
that the same data can produce different estimates because methods target
different populations and rely on different design choices.

\section{Common Mistakes}

\begin{boxF}
\textbf{Mistakes to avoid.}
\begin{itemize}
    \item Matching on post-treatment variables.
    \item Reporting a propensity-score regression table instead of balance.
    \item Estimating an ATE when there is overlap only for the treated
    population.
    \item Keeping extreme weights because the preferred estimate depends on
    them.
    \item Checking balance only on the propensity score and not on covariates.
    \item Forgetting that trimming, calipers, and dropped observations change
    the target population.
    \item Calling matching causal without defending conditional independence.
\end{itemize}
\end{boxF}

\section{Reporting a Matching or Propensity-Score Study}

A credible report should state:
\begin{enumerate}
    \item the causal estimand: ATE, ATT, ATU, or an overlap-population effect;
    \item the identifying assumption and why it is plausible in the setting;
    \item the list of covariates and why they are pre-treatment confounders;
    \item the matching method or propensity-score model;
    \item overlap diagnostics;
    \item balance before and after adjustment;
    \item number of observations dropped, matched, or reused;
    \item weight diagnostics if weighting is used;
    \item standard errors and how they account for the full procedure;
    \item sensitivity of results to calipers, trimming, covariates, and target
    estimand.
\end{enumerate}

The reader's main question is: after this design step, would I believe treated
and comparison units are exchangeable with respect to potential outcomes?

\section{Worked Problems}

\begin{exmp}[Exact matching ATT]
There are two cells. In cell A, treated mean outcome is 12 and control mean
outcome is 9. In cell B, treated mean outcome is 20 and control mean outcome is
18. Among treated units, 70 percent are in cell A and 30 percent are in cell B.
The ATT is
\[
    (12-9)\times 0.70 + (20-18)\times 0.30
    =
    2.1+0.6=2.7.
\]
The raw treated-control gap would generally be different because it also uses
the untreated distribution across cells.
\end{exmp}

\begin{exmp}[ATT weights]
Suppose a control unit has \(\widehat e(X)=0.80\). Its ATT control weight is
\[
    \frac{0.80}{1-0.80}=4.
\]
This control receives high weight because its covariates make it look like a
unit that was likely to be treated. If another control has
\(\widehat e(X)=0.20\), its ATT weight is \(0.20/0.80=0.25\). It contributes
less to the treated counterfactual.
\end{exmp}

\begin{exmp}[Effective sample size]
Consider weights \(w=(1,1,1,10)\). The nominal sample size is 4, but
\[
    ESS=\frac{(1+1+1+10)^2}{1^2+1^2+1^2+10^2}
    =
    \frac{169}{103}
    \approx 1.64.
\]
The weighted estimate is effectively using far less information than the raw
sample size suggests.
\end{exmp}

\section{Checklist for Students}

\begin{enumerate}
    \item Write the potential-outcomes estimand before choosing a method.
    \item State whether the target is ATE, ATT, ATU, or a trimmed/overlap
    population.
    \item Use only pre-treatment covariates for the design.
    \item Defend conditional independence using the institutional setting.
    \item Check overlap before trusting estimates.
    \item For matching, report the distance metric, replacement rule, number of
    neighbors, caliper, and dropped units.
    \item For propensity scores, report balance and overlap, not only the score
    model.
    \item For IPW, report weight percentiles, maximum weights, and effective
    sample size.
    \item Check standardized mean differences before and after adjustment.
    \item Explain how trimming or matching changes the population being studied.
\end{enumerate}

\section{Practice Questions}

\begin{enumerate}
    \item In your own words, explain why the raw treated-control difference can
    be biased in an observational study.
    \item Derive the exact-matching formula for the ATT using iterated
    expectations and conditional independence.
    \item Give an example of a post-treatment variable and explain why matching
    on it can be harmful.
    \item Why does exact matching become infeasible when many covariates are
    continuous?
    \item In nearest-neighbor matching, what is the tradeoff between one
    neighbor and many neighbors?
    \item What does overlap mean for the ATT? What does it mean for the ATE?
    \item State the propensity-score theorem and explain why it is useful.
    \item Why is high treatment-prediction accuracy not the main goal of a
    propensity-score model?
    \item Derive the IPW identity
    \(\E[D_iY_i/e(X_i)]=\E[Y_{1i}]\).
    \item A control unit has \(\widehat e(X)=0.90\). What is its ATE control
    weight? What is its ATT control weight? What does this imply about
    overlap?
    \item Run the companion Python script. Which estimator targets the ATT?
    Which targets the ATE? How do trimming and weights change the interpretation?
    \item Suppose a paper reports propensity-score matching but no balance table.
    What should you ask the author to provide?
\end{enumerate}

\end{document}
